\chapter*{Abreviaturas y siglas}
\addcontentsline{toc}{chapter}{Abreviaturas y siglas}

\begin{description}
  \item[AI4Mars] Conjunto de datos abierto de terreno marciano etiquetado por
    anotación colaborativa, publicado por el Jet Propulsion Laboratory.
  \item[APA] \textit{American Psychological Association}; estilo de citación
    empleado en este documento.
  \item[CSV] \textit{Comma-Separated Values}; formato de archivo tabular en texto
    plano utilizado para el conjunto de resultados.
  \item[EDR] \textit{Experiment Data Record}; nivel de procesamiento de los
    productos de imagen del Planetary Data System.
  \item[GEO] Taxonomía geológica extendida de AI4Mars, distinta de la escala de
    navegación y fuera del alcance de este trabajo.
  \item[HiRISE] \textit{High Resolution Imaging Science Experiment}; cámara orbital
    del \textit{Mars Reconnaissance Orbiter}.
  \item[IoU] \textit{Intersection over Union}; índice de solapamiento entre una
    región predicha y una de referencia.
  \item[JPL] \textit{Jet Propulsion Laboratory}, NASA.
  \item[MAHLI] \textit{Mars Hand Lens Imager}; cámara de aproximación del rover
    \textit{Curiosity}.
  \item[MER] \textit{Mars Exploration Rovers}; misión de los rovers
    \textit{Spirit} y \textit{Opportunity}.
  \item[MPS] \textit{Metal Performance Shaders}; interfaz de aceleración por
    unidad gráfica en procesadores Apple Silicon.
  \item[MSL] \textit{Mars Science Laboratory}; misión del rover \textit{Curiosity}.
  \item[NAV] Escala de clases de navegación de AI4Mars: suelo, lecho rocoso, arena
    y roca grande.
  \item[NavCam] \textit{Navigation Camera}; par estéreo de cámaras de navegación
    en blanco y negro montadas en el mástil del rover.
  \item[PDS] \textit{Planetary Data System}; archivo público de datos planetarios
    de la NASA.
  \item[SAM] \textit{Segment Anything Model}; modelo fundacional de segmentación
    de propósito general.
\end{description}
\clearpage
